\chapter{Performance and optimization}
\label{ch:performance}

\begin{quote}
\emph{``Premature optimization is the root of all evil --- but so is premature pessimization.''}
--- after Donald Knuth
\end{quote}

% ============================================================
\section{Why Julia is fast}

Julia achieves C-like speed through three pillars:

\begin{enumerate}
  \item \textbf{JIT compilation.} Julia compiles functions to native machine code via LLVM the first time they are called with specific argument types.
  \item \textbf{Type specialization.} \texttt{f(1, 2)} and \texttt{f(1.0, 2.0)} compile to different, optimized machine code.
  \item \textbf{Type inference.} The compiler traces types through function bodies, eliminating runtime type checks when all types are known.
\end{enumerate}

\begin{minted}{julia}
function mysum(v)
    s = zero(eltype(v))
    for x in v; s += x; end
    return s
end
mysum([1, 2, 3])        # compiles mysum(::Vector{Int64})
mysum([1.0, 2.0, 3.0])  # compiles mysum(::Vector{Float64})
\end{minted}

\begin{juliatip}
The first call to a function is slow (compilation latency). Julia 1.9+ uses package precompilation to reduce this. Use \texttt{PrecompileTools.jl} in your own packages.
\end{juliatip}

% ============================================================
\section{Measuring performance}

\begin{minted}{julia}
# @time — run twice! First call includes compilation time.
@time sum(rand(10^7))   # compilation
@time sum(rand(10^7))   # true timing

# BenchmarkTools.jl — accurate benchmarks for fast functions
using BenchmarkTools
v = rand(10_000)

@btime sum($v)          # runs many times, reports minimum
#   1.256 μs (0 allocations: 0 bytes)

@benchmark sum($v)      # full statistical report

# IMPORTANT: interpolate variables with $ to avoid global access cost
A = rand(100, 100)
@btime $A * $A          # correct
@btime A * A            # wrong — includes global lookup overhead
\end{minted}

\begin{warningbox}
Always interpolate input variables with \texttt{\$} in \texttt{@btime} and \texttt{@benchmark}. Without interpolation, the benchmark measures global variable access overhead, which skews results.
\end{warningbox}

% ============================================================
\section{Type stability}

A function is \emph{type-stable} if the return type depends only on argument types, not values:

\begin{minted}{julia}
# Type-UNSTABLE: return type depends on the VALUE of x
function unstable_sqrt(x)
    x >= 0 ? sqrt(x) : "negative"  # returns Float64 OR String!
end

# Type-STABLE: return type always matches input type
stable_sqrt(x) = x >= 0 ? sqrt(x) : NaN

# Common instability: changing variable type in a loop
function bad_sum(n)
    s = 0            # Int, becomes Float64 after s += i/2
    for i in 1:n; s += i / 2; end
    return s
end

function good_sum(n)
    s = 0.0          # Float64 from the start
    for i in 1:n; s += i / 2; end
    return s
end
\end{minted}

\begin{performancetip}
Type stability is the single most important factor for Julia performance. When a function is type-stable, the compiler generates code with no boxing, no dynamic dispatch, and no heap allocations for local variables.
\end{performancetip}

% ============================================================
\section{Inspecting compiled code}

\begin{minted}{julia}
f(x, y) = x^2 + y^2

@code_warntype f(1.0, 2.0)   # MOST USEFUL: red=instability, blue=stable
@code_typed f(1.0, 2.0)      # typed intermediate representation
@code_llvm f(1.0, 2.0)       # LLVM IR
@code_native f(1.0, 2.0)     # machine assembly
\end{minted}

\begin{minted}{julia}
# Detecting and fixing instability
unstable(x) = x > 0 ? x : 0        # returns Union{Float64, Int64}
stable(x)   = x > 0 ? x : zero(x)  # always returns typeof(x)

@code_warntype unstable(1.0)  # Body::Union{Float64, Int64} — RED
@code_warntype stable(1.0)    # Body::Float64 — BLUE
\end{minted}

% ============================================================
\section{Common performance pitfalls}

\subsection{Pitfall 1: Global variables}

\begin{minted}{julia}
x_global = 1.0   # BAD: type could change at any time

function bad_global()
    s = 0.0
    for i in 1:1000; s += x_global * i; end
    return s
end

good_local(x) = (s = 0.0; for i in 1:1000; s += x * i; end; s)

const X_CONST = 1.0   # GOOD: const lets the compiler optimize
good_const() = (s = 0.0; for i in 1:1000; s += X_CONST * i; end; s)

using BenchmarkTools
@btime bad_global()     # ~10 μs (slow, allocations)
@btime good_local(1.0)  # ~0.5 μs (20x faster)
\end{minted}

\subsection{Pitfall 2: Abstract containers and untyped structs}

\begin{minted}{julia}
v_any   = Any[1.0, 2.0, 3.0]       # BAD: boxed elements
v_typed = Float64[1.0, 2.0, 3.0]   # GOOD: contiguous memory

struct BadPoint; x; y; end                    # BAD: fields are ::Any
struct GoodPoint; x::Float64; y::Float64; end # GOOD: concrete types
struct FlexPoint{T<:Real}; x::T; y::T; end   # BEST: generic AND fast
\end{minted}

\subsection{Pitfall 3: Growing arrays in a loop}

\begin{minted}{julia}
# BAD: push! causes repeated reallocations
bad_build(n) = (v = Float64[]; for i in 1:n; push!(v, sin(i)); end; v)

# GOOD: preallocate
function good_build(n)
    v = Vector{Float64}(undef, n)
    for i in 1:n; v[i] = sin(i); end
    return v
end

# ALSO GOOD: comprehension (Julia preallocates internally)
good_build2(n) = [sin(i) for i in 1:n]
\end{minted}

% ============================================================
\section{Memory allocation and profiling}

\begin{minted}{julia}
# @allocated — returns bytes allocated
bytes = @allocated rand(1000, 1000)
println("Allocated $(bytes / 1024^2) MiB")

# In-place operations to reduce allocations
A = rand(1000, 1000)
A .= rand.(Float64)   # reuse memory, near-zero new allocations
\end{minted}

\begin{minted}{julia}
# Statistical profiler
using Profile

function heavy_work()
    A = rand(1000, 1000); B = A * A'; return sum(eigen(B).values)
end

heavy_work()               # warm up
@profile heavy_work()
Profile.print(format=:flat)
Profile.clear()

# Visual flame graph
using ProfileView
@profview heavy_work()     # wide bars = bottlenecks
\end{minted}

% ============================================================
\section{Practical: optimize a slow function step by step}

\begin{minted}{julia}
using BenchmarkTools, LinearAlgebra

# v1: naive — Vector{Any}, type instability
data = Any[rand() for _ in 1:100_000]
function slow_norm(v)
    s = 0  # Int!
    for i in 1:length(v); s = s + v[i]^2; end
    return sqrt(s)
end
@btime slow_norm($data)   # ~2.5 ms

# v2: fix container type
data_typed = rand(100_000)
@btime slow_norm($data_typed)   # ~1.2 ms (2x)

# v3: fix type instability
function better_norm(v)
    s = 0.0
    for i in eachindex(v); s += v[i]^2; end
    return sqrt(s)
end
@btime better_norm($data_typed)  # ~50 μs (50x)

# v4: @simd + @inbounds
function fast_norm(v)
    s = 0.0
    @inbounds @simd for i in eachindex(v); s += v[i]^2; end
    return sqrt(s)
end
@btime fast_norm($data_typed)    # ~15 μs (170x)

# v5: built-in BLAS
@btime norm($data_typed)         # ~12 μs (200x vs v1)
\end{minted}

\begin{performancetip}
The optimization checklist: (1) avoid globals, (2) concrete container types, (3) type stability, (4) preallocate, (5) \texttt{@views}, (6) \texttt{@inbounds @simd}, (7) \texttt{@threads} for parallelism.
\end{performancetip}

\begin{exercise}
\begin{enumerate}
  \item Write \texttt{my\_dot(a, b)} computing the dot product. Benchmark against \texttt{LinearAlgebra.dot} and optimize until within 2x.
  \item Use \texttt{@code\_warntype} to find and fix the instability in: \texttt{mystery(x) = (x > 0 ? x : 0) * 2.0}.
  \item Benchmark matrix-vector multiplication with \texttt{Vector\{Any\}} vs.\ \texttt{Vector\{Float64\}} for $n=1000$.
  \item Profile \texttt{heavy\_work()} and identify which sub-operation takes the most time.
  \item Compute $\sum_{k=1}^{n} \sin(k)/k$ via (a) loop, (b) generator, (c) broadcasting. Benchmark all three.
  \item Take a slow function from a previous exercise and apply the optimization checklist. Document each step.
\end{enumerate}
\end{exercise}

% ============================================================
\section{Chapter summary}

\begin{itemize}
  \item Julia achieves C-like speed through JIT compilation, type specialization, and type inference via LLVM.
  \item Use \texttt{@btime} from BenchmarkTools.jl for accurate benchmarks; always interpolate inputs with \texttt{\$}.
  \item Type stability is the most important performance factor: use \texttt{@code\_warntype} to detect instabilities.
  \item Avoid global variables, abstract containers (\texttt{Vector\{Any\}}), and untyped struct fields.
  \item Preallocate arrays and use \texttt{@views} to reduce memory allocations.
  \item Profile with \texttt{@profile} and visualize with ProfileView.jl to find bottlenecks.
\end{itemize}
